\par\Needspace{13\baselineskip}
\originalchapter{习题解答}
本章与前面各章的习题编号用链接相互对应. 各题为本书根据所讲概念另设或整理的标准训练题, 不宣称来自未核实的 MIT 作业题号.
\par\Needspace{9\baselineskip}
\originalsection{弧长与平面曲率}
\par\Needspace{5\baselineskip}
\noindent\textbf{习题\ref{ex:1a}.}
\begin{solution}
微分得 $\gamma'=e^t(\cos t-\sin t,\sin t+\cos t)$, $\gamma''=2e^t(-\sin t,\cos t)$. 所以 $v=\sqrt2e^t$, 行列式为 $2e^{2t}$, 代入曲率公式得 $k=e^{-t}/\sqrt2$. 在 $t=0$ 有 $\gamma=(1,0)$, $T=(1,1)/\sqrt2$, $n=(-1,1)/\sqrt2$, $k=1/\sqrt2$. 圆心为 $\gamma+n/k=(0,1)$, 半径为 $1/k=\sqrt2$.
\end{solution}
\par\Needspace{5\baselineskip}
\noindent\textbf{习题\ref{ex:1b}.}
\begin{solution}
利用 $n'=-kT$, 得
\[
 C'=T+\frac{n'}k-\frac{k'}{k^2}n=-\frac{k'}{k^2}n.
\]
因 $k\ne0$, 圆心轨迹正则当且仅当 $k'\ne0$. 在曲率临界点它的速度为0, 所给参数下失去正则性; 这不自动保证该点是尖点, 还需更高阶导数与非退化条件. 圆的 $k$ 恒定, 圆心轨迹是一点, 是最直接的例外.
\end{solution}
\par\Needspace{5\baselineskip}
\noindent\textbf{习题\ref{ex:1c}.}
\begin{solution}
取 $F=x^2+2y^2$, 梯度 $(2x,4y)$, Hessian 为 $\operatorname{diag}(2,4)$. 在 $(1,0)$, $J\nabla F=(0,2)$, Hessian 二次型为16, 梯度长度三次方为8, 得 $k=2$. 在 $(0,1/\sqrt2)$, $J\nabla F=(-2\sqrt2,0)$, 二次型为16, 分母为 $16\sqrt2$, 得 $k=1/\sqrt2$. 椭圆半轴 $a=1,b=1/\sqrt2$, 前章公式给出 $a/b^2=2$ 与 $b/a^2=1/\sqrt2$, 一致.
\end{solution}
\par\Needspace{9\baselineskip}
\originalsection{闭曲线}
\par\Needspace{5\baselineskip}
\noindent\textbf{习题\ref{ex:2a}.}
\begin{solution}
相对于指定周期, 它不能简单, 因简单正向曲线的总有向曲率只能为 $2\pi$, 反向为 $-2\pi$. 例子取 $\alpha(s)=R(\cos(s/R),\sin(s/R))$, 指定周期 $L=4\pi R$. 它为单位速度, $k=1/R>0$, 此周期走两圈, 总曲率为 $4\pi$. 该映射也有较小周期 $2\pi R$, 若用较小周期定义闭曲线, 则它是简单圆. 这说明必须固定所研究的一周.
\end{solution}
\par\Needspace{5\baselineskip}
\noindent\textbf{习题\ref{ex:2b}.}
\begin{solution}
令 $n_o=(\cos\varphi,\sin\varphi)$, $t=Jn_o$, 用 $\alpha=hn_o+h't$ 定义. 有 $h'=-2b\sin2\varphi$, $h''=-4b\cos2\varphi$, 所以
\[
 \rho=h+h''=a-3b\cos2\varphi>0,\qquad
 k=\frac1{a-3b\cos2\varphi}.
\]
曲线闭合且 $\alpha'=\rho t\ne0$. 为不循环地证明简单, 固定 $\varphi_0$ 并考察高度 $f(\varphi)=\langle\alpha(\varphi),n_o(\varphi_0)\rangle$. 其导数为 $\rho\sin(\varphi_0-\varphi)$, 在 $(\varphi_0,\varphi_0+\pi)$ 为负, 在下一半周为正. 因而一周内唯一最大点是 $\varphi_0$. 若另一个参数对应同一点, 它也达到该高度最大值, 矛盾. 于是曲线简单.

由 $\dd s=\rho\dd\varphi$, 有 $\dd k/\dd s=-6b\sin2\varphi/\rho^3$. $b\ne0$, 临界参数恰为 $0,\pi/2,\pi,3\pi/2$, 正好四个.
\end{solution}
\par\Needspace{5\baselineskip}
\noindent\textbf{习题\ref{ex:2c}.}
\begin{solution}
周长为 $\int\rho\dd\varphi=\int(h+h'')\dd\varphi=\int h\dd\varphi$, 最后一项由周期性 $\int h''=0$. 面积由 Green 公式为 $\frac12\int\det(\alpha,\alpha')\dd\varphi$. 因 $\alpha=hn_o+h't$, $\alpha'=\rho t$, $\det(n_o,t)=1$, 得行列式 $h\rho$. 因而
\[
 A=\tfrac12\int h(h+h'')\dd\varphi
 =\tfrac12\int(h^2-h'^2)\dd\varphi,
\]
分部积分的边界项因周期而为0. 曲线正向简单使此有向面积等于区域通常面积.
\end{solution}
\par\Needspace{9\baselineskip}
\originalsection{空间曲线}
\par\Needspace{5\baselineskip}
\noindent\textbf{习题\ref{ex:3a}.}
\begin{solution}
$\gamma'=(1,2t,3t^2)$, $\gamma''=(0,2,6t)$, $\gamma'''=(0,0,6)$. 叉积为 $(6t^2,-6t,2)$, 平方长度为 $4(1+9t^2+9t^4)$, 三重积为12. 因而
\[
 \kappa=\frac{2\sqrt{1+9t^2+9t^4}}{(1+4t^2+9t^4)^{3/2}},\qquad
 \tau=\frac3{1+9t^2+9t^4}.
\]
曲率处处正, 挠率处处正, 由平面性判据, 任意非空开区间都不可能落在一个固定平面上.
\end{solution}
\par\Needspace{5\baselineskip}
\noindent\textbf{习题\ref{ex:3b}.}
\begin{solution}
记 $r=\alpha-p$. $|r|^2=R^2$ 一次微分给 $\langle r,T\rangle=0$, 二次微分给 $1+\kappa\langle r,n\rangle=0$, 即 $\langle r,n\rangle=-1/\kappa$. 因而 $r=-n/\kappa+cb$. 再微分 $\langle r,n\rangle=-1/\kappa$, 得 $\langle r,n'\rangle=\kappa'/\kappa^2$. Frenet 方程使左侧为 $\tau c$, 所以在 $\tau\ne0$ 区间
\[
 r=-\frac1\kappa n+\frac{\kappa'}{\kappa^2\tau}b,\qquad
 R^2=\frac1{\kappa^2}+\left(\frac{\kappa'}{\kappa^2\tau}\right)^2.
\]
这是位于给定球面时的必要关系. 单独只知右侧恒定, 不能在此解答中未经进一步论证宣称已构造固定球心; 原题只要求从已在球面上推出该式.
\end{solution}
\par\Needspace{5\baselineskip}
\noindent\textbf{习题\ref{ex:3c}.}
\begin{solution}
设投影 $\beta=\alpha-\langle\alpha,U\rangle U$. 有 $\langle T,U\rangle=d/\sqrt{1+d^2}$, $U$ 恒定, 因而 $\beta'=T-dU/\sqrt{1+d^2}$, 速度 $v_\beta=1/\sqrt{1+d^2}>0$. 又 $\beta''=\kappa n$, 且 $n\perp U$, $n\perp\beta'$. 因速度为常数, 投影曲率为 $|\beta''|/v_\beta^2=(1+d^2)\kappa$. 它虽在平面上, 仍不必是圆, 因为 $\kappa$ 可随弧长变化.
\end{solution}
\par\Needspace{9\baselineskip}
\originalsection{曲面度量}
\par\Needspace{5\baselineskip}
\noindent\textbf{习题\ref{ex:4a}.}
\begin{solution}
$X_u=(1,0,f_u)$, $X_v=(0,1,f_v)$, 所以
\[
 \I=(1+f_u^2)\dd u^2+2f_uf_v\dd u\dd v+(1+f_v^2)\dd v^2,\qquad
 \dd A=\sqrt{1+f_u^2+f_v^2}\dd u\dd v.
\]
对 $f=(u^2+v^2)/2$, 极坐标下 $\dd A=\sqrt{1+r^2}r\dd r\dd\theta$. 面积为
\[
 2\pi\int_0^a r\sqrt{1+r^2}\dd r
 =\frac{2\pi}3\bigl((1+a^2)^{3/2}-1\bigr).
\]
积分区域在参数平面, 被积面积因子反映倾斜程度.
\end{solution}
\par\Needspace{5\baselineskip}
\noindent\textbf{习题\ref{ex:4b}.}
\begin{solution}
记 $q=u^2+v^2$, 分子范数平方为 $4q+(q-1)^2=(q+1)^2$, 故 $|X|=1$. 它不会到北极 $(0,0,1)$, 反之球面非北极点的逆为 $(u,v)=(x/(1-z),y/(1-z))$, 所以为单射光滑参数.

用商法则微分后可求得 $\langle X_u,X_u\rangle=\langle X_v,X_v\rangle=4/(1+q)^2$, $\langle X_u,X_v\rangle=0$. 于是 $\I=4(\dd u^2+\dd v^2)/(1+q)^2$, 正定且与欧氏度量成同一正因子, 所以保角. 它不保长度, 因伸缩因子随点变化.
\end{solution}
\par\Needspace{5\baselineskip}
\noindent\textbf{习题\ref{ex:4c}.}
\begin{solution}
有 $\I=(1+a^2)\dd r^2+r^2\dd\theta^2$. 令 $\rho=\sqrt{1+a^2}\,r$, $\psi=\theta/\sqrt{1+a^2}$, 则 $\I=\dd\rho^2+\rho^2\dd\psi^2$, 是平面极坐标度量. 在不绕过切缝的小邻域, $(\rho,\psi)\mapsto(\rho\cos\psi,\rho\sin\psi)$ 给出局部等距. 锥上一周 $\theta$ 增加 $2\pi$, 展开扇形角为 $2\pi/\sqrt{1+a^2}$. 顶点不在定义域, 局部可展并不能恢复顶点正则性.
\end{solution}
\par\Needspace{9\baselineskip}
\originalsection{曲面曲率}
\par\Needspace{5\baselineskip}
\noindent\textbf{习题\ref{ex:5a}.}
\begin{solution}
令 $r^2=x^2+y^2$, $D=1+r^2/a^2$. 图像的 Hessian 为 $a^{-1}I_2$, 所以 $K=1/(a^2D^2)$. 用平均曲率公式, 得
\[
 H=\frac{2+r^2/a^2}{2aD^{3/2}}.
\]
在 $r>0$ 点, 径向和环向为主方向, 相应主曲率为 $1/(aD^{3/2})$ 和 $1/(a\sqrt D)$. 这可用径向单位参数方向的第一基本形式值 $D$ 与第二基本形式值 $1/(a\sqrt D)$ 相除, 环向则第一基本形式值1而第二值相同. 两者处处正, 所有点为椭圆点. 只有 $D=1$, 即原点, 两主曲率相等, 是唯一脐点; 原点值都为 $1/a$.
\end{solution}
\par\Needspace{5\baselineskip}
\noindent\textbf{习题\ref{ex:5b}.}
\begin{solution}
Euler 公式给 $k_n=2\cos^2\theta-\sin^2\theta=3\cos^2\theta-1$. 零曲率方向满足 $\cos^2\theta=1/3$, 即 $\tan^2\theta=2$, 是两条无向直线. 最大值为2, 在第一主方向; 最小值为 $-1$, 在第二主方向. 对所有角作平均, $\cos^2$ 与 $\sin^2$ 的均值均为 $1/2$, 所以平均法曲率为 $(2-1)/2=1/2=H$.
\end{solution}
\par\Needspace{5\baselineskip}
\noindent\textbf{习题\ref{ex:5c}.}
\begin{solution}
对 $N=\nabla F/|\nabla F|$ 求导, 与切向 $W$ 取内积时归一化因子的导数项沿 $\nabla F$, 被切向正交消去. 所以 $\langle\dd N(V),W\rangle=\operatorname{Hess}F(V,W)/|\nabla F|$, 取负得到 $\II$.

在椭球点 $(a,0,0)$, 外法向为 $(1,0,0)$, 切平面由 $e_y,e_z$ 正交张成. $|\nabla F|=2/a$, Hessian 为 $\operatorname{diag}(2/a^2,2/b^2,2/c^2)$. 限制切平面得主曲率 $-a/b^2$ 与 $-a/c^2$, 分别沿 $y,z$ 方向; $K=a^2/(b^2c^2)>0$. 若改内法向, 两者同变号.
\end{solution}
\par\Needspace{9\baselineskip}
\originalsection{旋转与极小曲面}
\par\Needspace{5\baselineskip}
\noindent\textbf{习题\ref{ex:6a}.}
\begin{solution}
管外法向下 $k_\varphi=-1/a$, $k_\theta=-\cos\varphi/(R+a\cos\varphi)$, 故
\[
 H=-\frac{R+2a\cos\varphi}{2a(R+a\cos\varphi)}.
\]
若 $R=2a$, 令 $H=0$ 得 $\cos\varphi=-1$, 即内赤道 $\varphi=\pi$. 其余点 $H\ne0$, 所以环面不为极小曲面. 某一条曲线上平均曲率为0只是点集上的现象.
\end{solution}
\par\Needspace{5\baselineskip}
\noindent\textbf{习题\ref{ex:6b}.}
\begin{solution}
$X_r\times X_\theta=r(-a\cos\theta,-a\sin\theta,1)$, 法向为此向量除以 $r\sqrt{1+a^2}$. 第一形式系数 $E=1+a^2,F=0,G=r^2$; 第二形式 $e=f=0,g_2=ar/\sqrt{1+a^2}$. 故径向主曲率为0, 环向为 $a/(r\sqrt{1+a^2})$, 平均曲率为 $a/(2r\sqrt{1+a^2})$, 而 $K=0$. 这说明零 Gauss 曲率不等于极小.
\end{solution}
\par\Needspace{5\baselineskip}
\noindent\textbf{习题\ref{ex:6c}.}
\begin{solution}
取原点在曲面之外, 由紧致性取 $|p|=R>0$ 最大点. 在该点所有切向均垂直 $p$, 选法向 $N=p/R$. 对单位切向 $V$ 取通过点的曲面曲线, 距离平方二阶导数非正给出 $1+R\II(V,V)\le0$. 两个主方向上因此均有 $k_i\le-1/R$, 故 $H\le-1/R<0$. 若曲面极小则 $H=0$, 矛盾. 此论证不需要先选全局法向, 因 $H=0$ 对法向反转不变.
\end{solution}
\par\Needspace{9\baselineskip}
\originalsection{内蕴联络}
\par\Needspace{5\baselineskip}
\noindent\textbf{习题\ref{ex:7a}.}
\begin{solution}
度量矩阵为 $\operatorname{diag}(1,r^2)$, 只对 $u$ 变化. 代入 Christoffel 公式得
\[
 \Gamma^u_{vv}=-rr',\qquad \Gamma^v_{uv}=\Gamma^v_{vu}=r'/r,
\]
其余为0. 正交标架 $e_1=\partial_u,e_2=r^{-1}\partial_v$ 的连接形式为 $\omega=r'\dd v$, 外微分为 $r''\dd u\wedge\dd v$. 由 $\dd\omega=-K r\dd u\wedge\dd v$ 得 $K=-r''/r$. 同一个公式适用于不先给旋转嵌入的抽象度量.
\end{solution}
\par\Needspace{5\baselineskip}
\noindent\textbf{习题\ref{ex:7b}.}
\begin{solution}
$f_u=2u,f_v=2v$, 故 $\omega=-2v\dd u+2u\dd v$. 其外微分为 $4\dd u\wedge\dd v$, 面积元为 $e^{2(u^2+v^2)}\dd u\dd v$, 所以 $K=-4e^{-2(u^2+v^2)}<0$. 平面 Gauss 曲率为0, 局部等距必须保持曲率, 因而这个度量任何点附近都不能等距于欧氏平面.
\end{solution}
\par\Needspace{5\baselineskip}
\noindent\textbf{习题\ref{ex:7c}.}
\begin{solution}
取原正向单位标架及连接形式 $\omega$. 因 $K=0$, $\dd\omega=0$. 在足够小的星形坐标邻域, 闭一形式可积分为一个函数: 若 $\omega=A\dd u+B\dd v$, 从固定底边先横后竖积分, 用 $B_u=A_v$ 核查两个偏导即得到势函数. 取 $\psi$ 满足 $\dd\psi=-\omega$, 将标架旋转角 $\psi$, 新连接形式为0, 所以两标架向量 $\widetilde e_i$ 对所有方向都平行.

定义其对偶一形式 $\eta_i(V)=\langle\widetilde e_i,V\rangle$. 度量相容及无挠性给出
\[
 \dd\eta_i(V,W)=\langle\nabla_V\widetilde e_i,W\rangle-
 \langle\nabla_W\widetilde e_i,V\rangle=0.
\]
再用同样的局部积分法得到函数 $y_i$ 使 $\dd y_i=\eta_i$. 两一形式独立, 逆函数定理使 $(y_1,y_2)$ 成为局部坐标. 任意切向量的平方长度为 $\eta_1(V)^2+\eta_2(V)^2$, 故度量为 $\dd y_1^2+\dd y_2^2$, 映射到欧氏平面是局部等距. 单连通小邻域避免闭形式的周期障碍; 没有由此声称整张平直曲面全局等距于平面.
\end{solution}
\par\Needspace{9\baselineskip}
\originalsection{测地线}
\par\Needspace{5\baselineskip}
\noindent\textbf{习题\ref{ex:8a}.}
\begin{solution}
圆柱路径的水平单位方向可由引理\ref{lem:angle}取连续角函数, 故可以沿整条路径取展开坐标. 展开覆盖把终点提升为 $(R(\theta+2\pi m),h)$, $m\in\Z$. 从起点 $(0,0)$ 到每个提升点的直线投影给出一条测地线, 其长度为
\[
 L_m=\sqrt{R^2(\theta+2\pi m)^2+h^2}.
\]
反之每条圆柱测地线提升后是直线, 所以这些列全. 长度最小要求 $|\theta+2\pi m|$ 最小. 当 $0\le\theta<\pi$, 取 $m=0$; 当 $\pi<\theta<2\pi$, 取 $m=-1$; 当 $\theta=\pi$, 两者相同, 恰有两条不同最短轨迹. $\theta=0,h=0$ 为同点情形, 最短是常值路径, 非平凡绕圈都较长. 任意曲面路径提升后的欧氏长度不小于相应端点距离, 因而该比较也排除了非测地竞争路径.
\end{solution}
\par\Needspace{5\baselineskip}
\noindent\textbf{习题\ref{ex:8b}.}
\begin{solution}
由 $\dot\theta=c/r^2$, $\dot s=\pm\sqrt{1-c^2/r^2}$, 在 $\dot s\ne0$ 的区间有
\[
 \frac{\dd\theta}{\dd s}=\frac{\pm c}{r\sqrt{r^2-c^2}},
\]
符号与 $\dot s$ 相同. 在 $r=|c|$ 只知道经向速度瞬时为0, 可能是转向点. 要使整个纬线 $s=s_0$ 成为测地线, 必须由经向方程得到 $r'(s_0)=0$; 此条件也充分. 若 $r'\ne0$, $\ddot s=rr'\dot\theta^2$ 非零, 路径立即离开该纬线.
\end{solution}
\par\Needspace{5\baselineskip}
\noindent\textbf{习题\ref{ex:8c}.}
\begin{solution}
记反射限制为等距 $\sigma:M\to M$. 选固定曲线一点 $p$ 和单位切向 $V$. 因 $\sigma$ 在该曲线上恒等, $\sigma(p)=p,\dd\sigma_pV=V$. 具有初值 $p,V$ 的测地线 $\gamma$ 经等距后仍为测地线; 这来自能量第一变分或度量联络保持. $\sigma\circ\gamma$ 有相同初值, 唯一性给出 $\sigma\circ\gamma=\gamma$. 它因此留在局部固定点集.

题设固定点集局部是一条正则曲线, 所以此非恒定测地线在短时间内正好沿该曲线行进. 采用相同方向弧长参数的固定曲线与它相同, 因而为测地线. 对整曲线逐点应用即可. 正则固定曲线条件排除固定点集只有一点、具有分叉或为整个二维曲面的情形; 后一种情形不能据此把任意曲线都判为测地线.
\end{solution}
\par\Needspace{9\baselineskip}
\originalsection{距离与正常坐标}
\par\Needspace{5\baselineskip}
\noindent\textbf{习题\ref{ex:9a}.}
\begin{solution}
球面为闭连通嵌入曲面, 定理\ref{thm:closed-geodesic}给出最短测地线. 前章已求出所有非恒定球面测地线为大圆. 长度为 $\ell$ 的单位速度段满足 $\langle p,q\rangle/R^2=\cos(\ell/R)=\cos\alpha$. 非负长度中最小解为 $R\alpha$, 而大圆较短弧确实达到, 所以距离为此值.

若 $0<\alpha<\pi$, 两点与球心确定唯一过心平面, 较短大圆弧轨迹唯一. 若 $\alpha=\pi$, 任一包含直径 $p,-p$ 的过心平面给出两个半圆弧, 总共有无穷多条最短轨迹. 若 $\alpha=0$, 最短长度为0, 只有常值路径; 非恒定大圆绕圈不最短. 这些唯一性是轨迹及单调参数意义, 同一轨迹允许不同恒速区间表示.
\end{solution}
\par\Needspace{5\baselineskip}
\noindent\textbf{习题\ref{ex:9b}.}
\begin{solution}
任意路径长度至少为位移长度2. 对 $0<\eps<1$, 先从 $(-1,0)$ 沿实轴到 $(-\eps,0)$, 再走上半圆半径 $\eps$, 最后从 $(\eps,0)$ 走到 $(1,0)$, 全部留在穿孔平面. 总长为 $2(1-\eps)+\pi\eps\to2$, 故下确界为2.

若有路径达到2, 欧氏长度不等式的取等要求其速度几乎处处为非负水平向量; 对分段光滑路径可逐段得到纵坐标恒为0且横坐标单调. 连续性迫使它经过原点, 与定义域矛盾. 即使允许一般可求长路径, 长度取等也使每个中间点落在线段上并按顺序前进, 同样经过原点.
\end{solution}
\par\Needspace{5\baselineskip}
\noindent\textbf{习题\ref{ex:9c}.}
\begin{solution}
光滑指数映射满足 $F(r,\theta)=p+r e(\theta)+O(r^2)$, 且可对 $\theta$ 微分, 故 $F_\theta=r e'(\theta)+O(r^2)$, $|e'|=1$. 所以 $A(0,\theta)=0$, $A_r(0,\theta)=1$; 更具体地, $F_\theta=rB(r,\theta)$, $B$ 光滑且 $B(0,\theta)=e'(\theta)\ne0$, 所以 $A=r|B|$ 在 $r\ge0$ 有单侧光滑延拓, 这些初值可用于积分微分方程.

对 $r>0$, 第7章的正交度量公式给出 $A_{rr}=-K(F(r,\theta))A$. 两次积分及初值得
\[
 A(r,\theta)=r-\int_0^r(r-t)K(F(t,\theta))A(t,\theta)\dd t.
\]
由初始展开 $A(t)=t+O(t^2)$ 和 $K$ 有界, 积分为 $O(r^3)$, 所以先得到 $A=r+O(r^3)$. 再用 $K(F(t,\theta))=K(p)+o(1)$, 有
\[
 A=r-K(p)\int_0^r(r-t)t\dd t+o(r^3)
  =r-\tfrac16K(p)r^3+o(r^3).
\]
当 $K>0$ 时, 小圆的环向尺度比欧氏小, 与球面 $A=R\sin(r/R)$ 一致; 负曲率时则比欧氏大.
\end{solution}
\par\Needspace{9\baselineskip}
\originalsection{平行移动与局部公式}
\par\Needspace{5\baselineskip}
\noindent\textbf{习题\ref{ex:10a}.}
\begin{solution}
保持曲线参数方向, $N\mapsto-N$ 使 $N\times T$ 反向, 因而 $k_g\mapsto-k_g$; $S,\II$ 变号使 $k_n\mapsto-k_n$; 两主曲率一起变号, 所以 $K$ 不变.

Gauss--Bonnet 的边界不是任取方向, 而是区域在定向切平面的左侧. 反转曲面定向后, 正向边界也要反向. 曲线反向使 $k_g$ 再变一次号, 两次相消, 所以按照新正向边界计算的测地曲率积分数值保持. 曲面内角不变, 以新曲面定向及新边界方向同时定义的外角也不变. Gauss 曲率乘普通正面积的积分保持, 公式因此一致.
\end{solution}
\par\Needspace{5\baselineskip}
\noindent\textbf{习题\ref{ex:10b}.}
\begin{solution}
过心平面的交线给出互相垂直的球面大圆, 三个顶点均为直角, 内角和 $3\pi/2$. 球面被这三平面分成八个全等区域, 三角形面积为 $4\pi R^2/8=\pi R^2/2$. 乘 $K=1/R^2$ 得曲率积分 $\pi/2$. 因而 $\pi+\int K=3\pi/2$, 等于角和. 三个外角均为 $\pi/2$, 其总和与曲率积分合起来为 $2\pi$.
\end{solution}
\par\Needspace{5\baselineskip}
\noindent\textbf{习题\ref{ex:10c}.}
\begin{solution}
帽区域面积为 $2\pi R^2(1-\cos\theta_0)$, 由 holonomy 公式, 净旋转角为 $2\pi(1-\cos\theta_0)$ 模 $2\pi$. 当 $\theta_0=\pi/2$, 实数积分为 $2\pi$, 净旋转在单位向量上等于0角; 一个整周转动的终点向量与初始向量相同. 非零曲率并不保证每一个回路都有非恒等 holonomy, 它保证足够小且有非零曲率积分的回路出现相应小转角.
\end{solution}
\par\Needspace{9\baselineskip}
\originalsection{整体公式}
\par\Needspace{5\baselineskip}
\noindent\textbf{习题\ref{ex:11a}.}
\begin{solution}
对连通曲面, $-\Area=\int K=4\pi(1-g)$, 故 $\Area=4\pi(g-1)$. 面积为正迫使 $g\ge2$. 这只是内蕴几何的必要拓扑与面积条件, 不自动产生三维欧氏嵌入; 事实上若它是紧致无边界嵌入曲面, 命题\ref{prop:elliptic-point}要求存在 $K>0$ 点, 与 $K=-1$ 矛盾. 抽象曲面与嵌入曲面的范围必须区分.
\end{solution}
\par\Needspace{5\baselineskip}
\noindent\textbf{习题\ref{ex:11b}.}
\begin{solution}
曲率积分为
\[
 \int K\dd A=\int_0^{2\pi}\int_{\theta_1}^{\theta_2}\sin\theta\dd\theta\dd\varphi
 =2\pi(\cos\theta_1-\cos\theta_2).
\]
在北边界 $\theta_1$, 区域位于更大的极角一侧, 正向边界沿 $\varphi$ 减小, 所以该边积分为 $-2\pi\cos\theta_1$. 南边界 $\theta_2$ 沿 $\varphi$ 增加, 积分为 $2\pi\cos\theta_2$. 两者之和恰抵消曲率积分. 纬带为环形区域, $\chi=0$, 无角点, 故完全满足整体公式.
\end{solution}
\par\Needspace{5\baselineskip}
\noindent\textbf{习题\ref{ex:11c}.}
\begin{solution}
平面曲率 $K=0$, 环形区域 $\chi=0$, 无角点, 所以两边界测地曲率积分总和为0. 若取向上法向, 外边界使区域在左, 按逆时针走, 是正向简单闭曲线, 转角定理给 $2\pi$. 内边界为使环形区域在左必须顺时针走, 积分为 $-2\pi$. 同一条内边界若单独作为其所围小圆盘的边界会取相反方向; 不能忽略其所界区域而只凭曲线形状定向.
\end{solution}
\par\Needspace{9\baselineskip}
\originalsection{双曲平面}
\par\Needspace{5\baselineskip}
\noindent\textbf{习题\ref{ex:12a}.}
\begin{solution}
由距离公式, $d(0,1/2)=2\arctanh(1/2)=\log3$. 两点 $-1/2,1/2$ 对应的 $q=1/(1+1/4)=4/5$, 所以 $d(-1/2,1/2)=\log9=2\log3$. 沿实轴, $x(s)=\tanh((s-s_0)/2)$ 有 $2|x'|/(1-x^2)=1$, 是单位速度参数, 可选 $s_0$ 调整起点. 方向反向取其负号或反向参数.
\end{solution}
\par\Needspace{5\baselineskip}
\noindent\textbf{习题\ref{ex:12b}.}
\begin{solution}
三角形角和为 $\pi(1/3+1/4+1/6)=3\pi/4$, 双曲面积为 $\pi/4$. 若单位球面中有这样的测地圆盘三角形, 则 $K=1$ 的 Gauss--Bonnet 会给面积等于角和减 $\pi=-\pi/4$, 不可能. 所以不存在所指球面三角形. 此判断按简单测地圆盘和实际内角作出, 不把跨多周或带重数的三条弧称为同一种三角形.
\end{solution}
\par\Needspace{5\baselineskip}
\noindent\textbf{习题\ref{ex:12c}.}
\begin{solution}
度量是 $e^{2f}(\dd x^2+\dd y^2)$, $f=-\log y$, $\Delta f=1/y^2$, 所以 $K=-1$. Christoffel 公式给出 $\Gamma^x_{xy}=\Gamma^x_{yx}=-1/y$, $\Gamma^y_{xx}=1/y$, $\Gamma^y_{yy}=-1/y$, 其余为0. 若 $x=x_0$, 测地线方程变为 $y''-y'^2/y=0$, 解 $y(t)=Ae^{ct}$, $A>0$. 速度为 $|c|$, 所以 $c=\pm1$ 是单位速度.

任意连接同一竖线两点的路径长度满足
\[
 L=\int\frac{\sqrt{x'^2+y'^2}}y\dd t
 \ge\int\frac{|y'|}y\dd t
 \ge|\log y_2-\log y_1|.
\]
竖直单调路径取等, 因此距离为 $|\log(y_2/y_1)|$. 第一不等式取等要求 $x'=0$, 第二要求高度不回走, 也给出最短轨迹唯一.
\end{solution}
